\section{Controller Design}
\label{sec:control}

\subsection{Closed-Loop Requirements}
\label{sec:requirements}
Both control laws are judged against the closed-loop requirements of
Table~\ref{tab:requirements}, adopted from the design practice of our
research group~\cite{mirko2026companion} and anchored, row by row, in
the flight-control literature. The stability-margin rows descend from the military
flight-control-system specification MIL-F-9490D, which requires a gain
margin of at least 6\,dB and a phase margin of at least $45^\circ$ for
the augmented airplane~\cite{milf9490d1975} --- limits carried into
present-day industry practice by its successor SAE
AS94900~\cite{sae2007as94900}; the 10\,dB adopted here is deliberately
conservative. The damping row is stricter than the flying-qualities
levels of MIL-F-8785C and MIL-STD-1797A, whose Level-1 bands admit
short-period and Dutch-roll damping well below
0.75~\cite{milf8785c1980,moorhouse1982milf8785c,milstd1797a1990}; the
$\zeta>0.75$ target reflects the classical well-damped second-order
response ($\zeta\approx0.7$ yields ${\approx}5\%$
overshoot~\cite{franklin2018feedback}). The time-domain rows
(settling time, overshoot, steady-state error) are outer-loop
performance targets in the sense of classical autopilot design
texts~\cite{blakelock1991automatic,nelson1998flight}, sized for this
platform's mission profile. Since these criteria were written for
piloted aircraft, their transfer to a small UAV is itself a design
decision, supported by the literature on scaled flying-qualities
requirements for unmanned aircraft~\cite{cotting2010evolution}.

\begin{table}[!t]
\caption{Closed-loop design requirements and their scope.}
\label{tab:requirements}
\centering
\begin{tabular}{lll}
\toprule
Requirement & Value & Scope \\
\midrule
Settling time & $<60$\,s & guidance channels ($\VT$, $\psi$, $h$) \\
Overshoot & $<10\%$ & guidance channels ($\VT$, $\psi$, $h$) \\
Steady-state error & $2$--$5\%$ & guidance channels ($\VT$, $\psi$, $h$) \\
Damping ratio & $\zeta>0.75$ & augmented inner modes \\
Gain margin & $>10$\,dB & each loop, at design \\
Phase margin & $>45^\circ$ & each loop, at design \\
\bottomrule
\end{tabular}
\end{table}

Each row carries an explicit scope. The time-domain rows apply to the
\emph{guidance-level} channels ($\VT$, $\psi$, $h$) on shaped
(prefiltered) mission references --- they are verified against both
controllers in Section~\ref{sec:sil-comparison}. The damping row
applies to the augmented inner modes at design time; the margin rows
are enforced per loop during loop shaping (the phase-margin targets of
Table~\ref{tab:gains}, $55$--$65^\circ$, exceed the requirement). One
honest caveat closes the picture: the augmented Dutch-roll damping
achieved by the PID design ($0.707$, Section~\ref{sec:pid}) satisfies
the MIL-F-8785C Level-1 band with wide margin but sits marginally
below the stricter internal target of $0.75$ --- reported as designed
rather than adjusted to fit. On the short-period side, the pitch-rate
damper gain of the tuning of record ($K_q=0.10$,
Section~\ref{sec:pid}) was raised precisely to lift the short-period
damping.

\subsection{Cascade PID Architecture}
\label{sec:pid}
The PID controller follows a classical cascade architecture ---
successive loop closure in the sense of~\cite{beard2012small} ---
designed loop by loop on the linearized models of
Section~\ref{sec:model} and then verified on the full nonlinear model.
The design methodology --- loop ordering, damper-first stability
augmentation, and the coordinated-turn heading law ---
follows~\cite{santos2018}, developed there for a conventional
general-aviation aircraft and applied here to the hybrid VTOL platform
in its fixed-wing flight mode.

\subsubsection{Stability Augmentation}
Rate dampers are closed first, tuned by root locus on the linear
subsystems. The pitch-rate damper $K_q\,q$ raises the short-period
damping ($K_q=0.10$ in the tuning of record, a deliberately raised
damping target); the yaw damper feeds $r$ through a washout
$s/(s+1)$ to the rudder, raising Dutch-roll damping to
$\zeta_{DR}\approx0.707$ without opposing steady turns. The washout
yaw damper is one of the oldest devices of automatic flight control:
introduced with the first swept-wing jets --- which could not be flown
unaugmented because of their Dutch roll --- in the late
1940s~\cite{mcruer1981eighty}, it received its definitive analytical
treatment in the classical stability-augmentation
literature~\cite{mcruer1973aircraft,blakelock1991automatic}, which the
present implementation follows. The roll
subsidence mode already meets Level~1 handling qualities with margin
($\lambda=-11.3$\,s$^{-1}$, Table~\ref{tab:poles}), so no roll-rate damper
is used ($K_p=0$).

\subsubsection{Longitudinal Loops}
\begin{itemize}
  \item \emph{Pitch-attitude inner loop} ($C_{\theta}$): PI on
        $\theta^{ref}-\theta$ to the elevator, closed over the $K_q$ damper;
  \item \emph{Airspeed hold} ($C_{vel}$): PI on $\VT^{ref}-\VT$ to the
        throttle;
  \item \emph{Altitude outer loop} ($C_{alt}$): PI on $h^{ref}-h$ to the
        pitch-attitude reference, with the command clamped to
        $\pm10^\circ$ and clamping-type anti-windup. The clamp is an
        angle-of-attack protection: with $\theta^{ref}$ limited to
        $+10^\circ$ about the trim, the peak $\alpha$ during an altitude
        capture stays near $16^\circ$ (the plant model has no stall
        representation, so the limit is imposed by design).
\end{itemize}
Depending on the flight mode, the pilot-level selector engages either the
altitude cascade (\emph{AltitudeHold}) or direct pitch tracking
(\emph{ThetaHold}), in which case the altitude loop is open.

\subsubsection{Lateral--Directional Loops}
\begin{itemize}
  \item \emph{Bank inner loop} ($C_{\phi}$): PI on $\phi^{ref}-\phi$ to
        the aileron (negative gains, since $C_{l_{\delta_a}}<0$);
  \item \emph{Heading hold} ($K_{heading}$): proportional law
        $\phi^{ref}=K_{heading}(\psi^{ref}-\psi)$ with
        $K_{heading}=u^e/(g\,\tau_\psi)$ and $\tau_\psi=6$\,s, the
        coordinated-turn relation~\cite{santos2018}. As with the
        longitudinal channel, the mode selector engages either heading
        hold (\emph{PsiHold}) or direct bank commands (\emph{PhiHold}).
\end{itemize}

\subsubsection{Tuning and Resulting Gains}
The compensators are PI ($K_d=0$): damping is provided by the explicit
dampers rather than by derivative action, which also avoids differentiating
noisy rate measurements on the embedded target. Each loop was tuned by
loop shaping (\texttt{pidtune}) with \emph{sequential inner-first} closure
--- every outer loop is designed against the already-closed inner
dynamics --- targeting the crossover frequencies and margins listed in
Table~\ref{tab:gains}, with a bandwidth separation of 5--7$\times$ between
consecutive loops.

\begin{table}[!t]
\caption{PID loops: design targets and resulting gains (tuning of
record).}
\label{tab:gains}
\centering
\begin{tabular}{llccl}
\toprule
Loop & Map & $K_P$ & $K_I$ & Original target \\
\midrule
$C_\theta$ & $\theta^{ref}\!\to\!\delta_e$ & 0.917 & 0.857 & $\omega_c$=2.0, PM=65$^\circ$ \\
$C_{vel}$ & $\VT^{ref}\!\to\!\delta_t$ & 0.32 & 0.096 & $\omega_c$=0.8 \\
$C_{alt}$ & $h^{ref}\!\to\!\theta^{ref}$ & 0.050 & 0.0043 & $\omega_c$=0.35, PM=55$^\circ$ \\
$C_\phi$ & $\phi^{ref}\!\to\!\delta_a$ & $-0.45$ & $-0.036$ & $\omega_c$=2.0, PM=65$^\circ$ \\
$K_{heading}$ & $\psi^{ref}\!\to\!\phi^{ref}$ & 0.1975 & --- & $\tau_\psi$=6\,s \\
$K_q$ & $q\!\to\!\delta_e$ & 0.10 & --- & $\zeta_{SP}$ raised \\
$K_r$ & $r\!\to\!\delta_r$ (washout) & 0.1481 & --- & $\zeta_{DR}$=0.707 \\
\bottomrule
\end{tabular}
\end{table}

The outer altitude and airspeed loops use two-degree-of-freedom
setpoint weighting with $b=0$ --- the proportional term acts on the
measured output only, with the reference entering through the integral
path --- and this tuning of record, retuned against the requirements
of Table~\ref{tab:requirements}, is the configuration used in every
campaign of this article, matching the public repository HEAD.

A command prefilter $1/(\tau_{ref}s+1)$ with $\tau_{ref}=3$\,s is applied
to the \emph{increment} of each reference before the loops, giving a
two-degree-of-freedom structure in the tradition of reference command
shaping~\cite{singer1990preshaping}: the command response is smoothed
--- the raw step would otherwise demand an unrealizable actuator jump
--- without altering the feedback design, its margins, or the clamp. This detail
matters for the embedded deployment: in the HIL architecture of
Section~\ref{sec:hil} the references travel through the link
\emph{already filtered}, so the prefilter stays on the host and only the
feedback path is embedded.

With all loops closed, the linear closed-loop system is stable with every
pole in the left half-plane: the slowest mode is at $-0.127$\,s$^{-1}$ and
the minimum damping ratio among the oscillatory modes is $\zeta=0.55$.
Notably, the open-loop unstable spiral mode ($+0.233$\,s$^{-1}$,
Table~\ref{tab:poles}) is stabilized at $-0.143$\,s$^{-1}$.

\subsection{Output-Feedback LQR (LQRy)}
\label{sec:lqry}
The second control law is a linear--quadratic design with \emph{output}
weighting (LQRy): for each loop, the quadratic cost
\begin{equation}
J=\int_0^\infty\!\big(y^{\!\top}Q\,y+u^{\!\top}R\,u\big)\,dt,\qquad y=C_{syn}x,
\end{equation}
is minimized over the linearized models of Section~\ref{sec:model},
yielding constant gains on the measured states plus integral action on
the tracked error. This is the classical linear--quadratic regulator
of Kwakernaak and Sivan~\cite{kwakernaak1972linear} in the
constant-gain \emph{output-feedback} setting introduced by Levine and
Athans~\cite{levine1970output} --- a problem whose difficulty and
solution landscape are surveyed in~\cite{syrmos1997static}, and whose
standard numerical solution is the flight-control-motivated iterative
algorithm of Moerder and Calise~\cite{moerder1985convergence}; the
formulation and its application to aircraft control follow the
treatment of Stevens \emph{et al.}~\cite{stevens2016aircraft}. The
weight-selection rationale follows established LQ design
practice~\cite{anderson1990optimal}, starting from Bryson's
rule~\cite{bryson1975applied}. The design compared in this article is
the \emph{gain-scheduled} LQRy of the companion article by the second
author~\cite{mirko2026companion}, which details its synthesis ---
including the target-transmission-zero augmentation and the gain
scheduling over the flight envelope --- and its frequency-domain
(Nichols) robustness assessment; it is summarized here.

The architecture deliberately mirrors the PID mode structure, so that
the two laws face the same tasks: five independent designs ---
\emph{ThetaHold}, \emph{AltitudeHold} (cascading into
$\theta^{ref}$), \emph{VelHold}, \emph{PhiHold} (a $2\times2$ design
commanding aileron and rudder), and \emph{PsiHold} (cascading into
$\phi^{ref}$) --- selected by the same three mode switches used by the
PID (Table~\ref{tab:lqry-loops}). Each synthesis model augments the
relevant linear subsystem with a frequency-selective filter
\begin{equation}
G(s)=\frac{s^{2}+2\zeta\omega_{n}s+\omega_{n}^{2}}{s}
\end{equation}
on each tracked output --- imposing target transmission zeros that
shape the closed loop while providing the integral action --- plus a
first-order actuator lag at 24\,rad/s; the constant output-feedback
gains are obtained from the coupled Lyapunov equations of the LQ
output-feedback problem~\cite{levine1970output,moerder1985convergence}.
The design is synthesized at nine envelope conditions
(12/15/18\,m/s $\times$ 590/600/610\,m); the SIL comparison of
Section~\ref{sec:sil} exercises the condition-2 (12\,m/s, 600\,m) gain
set, consistent with the mission's initial trim --- mid-mission gain
switching is not exercised. The weights per condition are given in the
companion article~\cite{mirko2026companion}.

\begin{table}[!t]
\caption{LQRy loops: regulated output and actuator (weights per
envelope condition in the companion article~\cite{mirko2026companion}).}
\label{tab:lqry-loops}
\centering
\begin{tabular}{lll}
\toprule
Loop & Output & Actuator \\
\midrule
ThetaHold & $\theta$ & elevator \\
AltitudeHold & $h\!\to\!\theta^{ref}$ & elevator \\
VelHold & $\VT$ & throttle \\
PhiHold & $\phi$, $\beta$ & ail.+rud. \\
PsiHold & $\psi\!\to\!\phi^{ref}$ & aileron \\
\bottomrule
\end{tabular}
\end{table}

Table~\ref{tab:lqry-gains} records the condition-2 gains: these are
the values used in the SIL comparison of Section~\ref{sec:sil}. Two
structural remarks: the pitch design feeds back the elevator
(actuator) state itself, its synthesis model carrying the 24\,rad/s
lag as a fifth state; and the $2\times2$ lateral design carries a
sideslip-integral gain alongside the bank integral (0.00845 on the
aileron and 0.02023 on the rudder at condition~2) --- small relative
to the $\phi$ integral, but not negligible. The earlier
\emph{fixed-gain} LQRy delivery, deployed on the HIL bench of
Section~\ref{sec:hil}, used a different, softer gain set, available in
the repository~\cite{repo}.

\begin{table}[!t]
\caption{LQRy gains at envelope condition 2 (12\,m/s, 600\,m), as used
in the SIL comparison (4 significant digits). Convention:
$u=G_{x}\,x+K_{I}\!\int\!e\,dt$. The $K_I$ column of PhiHold is the
bank ($\phi$) integral gain.}
\label{tab:lqry-gains}
\centering
\resizebox{\linewidth}{!}{%
\begin{tabular}{llll}
\toprule
Loop & State vector $x$ & $G_{x}$ & $K_I$ \\
\midrule
ThetaHold & $[\VT\,\alpha\,q\,\theta\,\delta_e]$ &
$[-0.00461\;6.477\;-15.60\;-483.2\;-6.795]$ & $-625$ \\
AltitudeHold & $[\VT\,\alpha\,q\,\theta\,h]$ &
$[-0.0718\;4.474\;-0.2121\;-5.064\;-0.5086]$ & $-0.2025$ \\
VelHold & $[\VT\,\alpha\,q\,\theta]$ &
$[-78.81\;-0.9055\;0.0589\;1.306]$ & $-50$ \\
PhiHold (ail.) & $[\beta\,p\,r\,\phi]$ &
$[0.00693\;0.01013\;0.03505\;0.1595]$ & $0.1720$ \\
PhiHold (rud.) & $[\beta\,p\,r\,\phi]$ &
$[-0.01085\;-0.02236\;-0.08167\;-0.3542]$ & $-0.3595$ \\
PsiHold & $[\beta\,p\,r\,\phi\,\psi]$ &
$[-5.387\;-0.1123\;-0.7948\;-1.272\;-6.037]$ & $-7.773$ \\
\bottomrule
\end{tabular}}
\end{table}

Unlike the PID harness, the LQRy harness of the gain-scheduled design
includes first-order actuator lags (24\,rad/s on the surfaces, 0.1\,s
on the throttle) and a $[0,1]$ throttle saturation, but no rate
limits, no surface position saturations, no reference prefilter, and
no $\theta^{ref}$ clamp --- references enter the loops as raw steps.
These differences between the two design harnesses are consolidated
and discussed in Section~\ref{sec:protocol}.

\subsection{Roll-to-Longitudinal Coupling}
\label{sec:coupling}
Both designs treat the longitudinal and lateral--directional channels
independently, so the classical roll-to-pitch coupling --- banking
tilts the lift vector, and the lift deficit, proportional to
$1-\cos\phi$, drops the nose and costs altitude --- is handled by
feedback alone. Since this is a deliberate design choice, it was
quantified rather than assumed benign: both controllers --- the PID
tuning of record and the gain-scheduled LQRy --- were subjected
to bank captures of increasing amplitude ($\phi^{ref}$ steps of
$5/15/30^\circ$ at $t=20$\,s, each rig shaping the reference as its
harness prescribes), in \emph{ThetaHold} (altitude loop open, the
worst case) and in \emph{AltitudeHold}, with all longitudinal metrics
taken as deviations from each rig's own $\phi=0$ baseline
(Fig.~\ref{fig:coupling}).
\begin{figure}[!t]\centering
\includegraphics[width=\linewidth]{fig_acoplamento.png}
\caption{Roll-to-longitudinal coupling, $\phi^{ref}=30^\circ$ step at
$t=20$\,s: $\phi$ (with the reference), $\Delta\theta$,
$\Delta h$, and $\Delta$elevator relative to each configuration's
$\phi=0$ baseline. Solid: ThetaHold at the 12\,m/s trim; dashed:
AltitudeHold at 15.2\,m/s.}\label{fig:coupling}\end{figure}

The results confirm the mechanism and bound its size. For the PID in
ThetaHold, the transient pitch dip scales with $1-\cos\phi$ as
expected --- $-0.005^\circ$, $-0.04^\circ$, and $-0.18^\circ$ for
$5/15/30^\circ$ --- negligible at the $5^\circ$ amplitude of the
flight campaign. The LQRy dip is $4$--$5\times$ larger
(${\approx}-0.8^\circ$ at $30^\circ$): its multivariable structure
does \emph{not} natively absorb the coupling better than the cascade.
With the altitude loop open the altitude drifts by construction ---
$-15.8$\,m for the PID over 200\,s at $30^\circ$, while the LQRy,
whose high-bandwidth airspeed loop adds energy through the turn,
drifts $+10.6$\,m \emph{upward}; closing
the altitude loop absorbs the coupling, with $\theta^{ref}$ settling
higher ($+0.59^\circ$ PID, $+0.71^\circ$ LQRy) and $\Delta h$
returning to zero after a transient dip of $0.10$\,m (PID) and
$1.45$\,m (LQRy), at the cost of a small steady
elevator increment.

In the steady turn, the elevator increment is \emph{larger} in the
ThetaHold runs at the 12\,m/s trim (solid traces, $2.06^\circ$ PID
and $2.26^\circ$ LQRy) than in the AltitudeHold runs at 15.2\,m/s
(dashed, $1.34^\circ$ and $1.55^\circ$), which at first seems
backwards ---
why would the altitude-holding configuration need \emph{less}
elevator? The explanation is the operating point, not the mode: sustaining the same $30^\circ$ bank
(load factor $1/\cos 30^\circ = 1.15$) requires an increment of lift
that, at 15.2\,m/s, is produced with less angle of attack --- and less
elevator --- because the dynamic pressure is $(15.2/12)^2 = 1.6$ times
higher. The measured increment ratios ($0.65$ for the PID, $0.69$ for
the LQRy) sit within a few
percent of the dynamic-pressure prediction $(12/15.2)^2 = 0.62$, the
residual owing to the $\theta^{ref}$ adjustment and the
angle-of-attack dependence of drag.

In answer to the natural question --- no, no feedforward decoupling
term is used by either law ($K_{f\!f}=0$ throughout): with the
altitude loop closed the residual coupling is $0.10$\,m even at a
$30^\circ$ bank (PID), and at the $5^\circ$ amplitude of the flight
campaign the $1-\cos\phi$ factor is a further $35\times$ smaller, so
feedback suffices.
A feedforward stage, and its evaluation on both laws at larger bank
angles, is left as future work
(Section~\ref{sec:conclusion}).
